% Generated by `taskgen course examples` from crypto.congruential (seed 3).
% Do not edit: change examples.yaml and regenerate.
\begin{example}[атака на конгруэнтную криптосистему]\label{ex:lattice-congruential}
Открытый текст $m$ зашифрован конгруэнтной криптосистемой. Зная шифротекст $c = 12002386$ и открытый ключ \[(q, h) = (21350387, 5352986),\] восстановить $m$.

\par\noindent\textit{Решение.} Закрытый ключ $(f, g)$ удовлетворяет $fh \equiv g\ (\textrm{mod}\ q)$ и мал, т.е. $(f, g)$ --- короткий вектор решётки с базисом $(1, h)^T = (1, 5352986)^T$, $(0, q)^T = (0, 21350387)^T$. Ищем кратчайший вектор методом Гаусса:\par Алгоритм Гаусса: если $\|v_2\| < \|v_1\|$ "--- меняем местами (${}^{\leftrightarrow}$); $m = \lfloor v_1 \cdot v_2 / \|v_1\|^2 \rceil$; если $m = 0$ "--- стоп, иначе $v_2 \gets v_2 - m v_1$.
\[\begin{array}{c|r|r|r}
 & v_1 & v_2 & m \\ \hline
1 & (1, 5352986) & (0, 21350387) & 4 \\
2^{\leftrightarrow} & (-4, -61557) & (1, 5352986) & -87 \\
3^{\leftrightarrow} & (-347, -2473) & (-4, -61557) & 24 \\
4 & (-347, -2473) & (8324, -2205) & 0 \\
\end{array}\]
Кратчайший вектор $\pm(-347, -2473)^T$; берём знак с $f' > 0$: $(f', g') = (347, 2473)$ (годится любой вектор, с которым расшифрование корректно). Расшифрование: $a \equiv f' c \equiv 347 \cdot 12002386 \equiv 1502477\ (\textrm{mod}\ q)$, $f'^{-1} \equiv 1967\ (\textrm{mod}\ |g'|)$, $b \equiv a f'^{-1} \equiv 1244\ (\textrm{mod}\ |g'|)$.\[m = 1244.\]
\end{example}
